% !PRINTED 136
\weberpage{136}
%<<CONT
ist dann in keiner der Gruppen $Q_1, Q_2, \dots, Q_{q-1}$ enthalten, weil angenommen war, dass $Q$ mit keiner dieser Gruppen einen Theiler gemein hat, dessen Grad grösser als $p^{\gamma}$ ist. Ebenso können wir einen Theiler $R_1$ von $Q_1$ vom Grade $p^{\gamma+1}$ finden, von dem $Q_0$ Normaltheiler ist, und der in keiner der Gruppen $Q, Q_2, \dots, Q_{q-1}$ enthalten ist. Nun ist $R$ sowohl als $R_1$ in $P$ enthalten. Wir betrachten die kleinste Gruppe $S$, die $R$ und $R_1$ zugleich enthält, die jedenfalls in $P$ enthalten ist (das kleinste gemeinschaftliche Vielfache von $R$ und $R_1$). Der Grad dieser Gruppe $S$ kann nicht eine Potenz von $p$ sein, da sonst nach dem Satze II., \S.\,30 $S$ und mithin $R$ und $R_1$ in einer der Gruppen $Q, Q_1, \dots, Q_{q-1}$ enthalten sein müssten. Es ist aber $R$ nur in $Q$, $R_1$ nur in $Q_1$ enthalten; also ist diese Annahme unzulässig. Der Grad von $S$ muss also von der Form $p^{\lambda} q$ sein.

Es kann aber $\lambda$ nicht kleiner als $\alpha$ sein; denn es ist, da eine Gruppe vom Grade $p^{\gamma+1}$ in $S$ enthalten ist, nämlich $R$,

\weq{4}{\lambda \geqq \gamma+1 > \beta.}

Nach der in $\alpha)$ gemachten Voraussetzung ist aber in $P$ kein Theiler vom Grade $p^{\lambda} q$ enthalten, in dem $\lambda$ zugleich grösser als $\beta$ und kleiner als $\alpha$ ist. Nach (4) ist also nothwendig $\lambda = \alpha$, d.\,h. $S$ ist mit $P$ identisch.

Nun ist $Q_0$ ein Normaltheiler von $R$ und $R_1$. Fassen wir also alle Elemente $c$ von $P$, die der Bedingung

\[ c^{-1}\,Q_0\,c = Q_0 \]

genügen, zu einer Gruppe zusammen, so enthält diese Gruppe sowohl $R$ als $R_1$, und ist also die ganze Gruppe $P$. Es ist also $Q_0$ auch Normaltheiler von $P$. Da nach (3) $\gamma$ grösser als Null ist, so ist der Grad von $Q_0$ grösser als 1, und wir stossen auf einen Widerspruch mit der Annahme, dass $P$ einfach sei.

\begin{center}
\S.\,34.

\gesperrt{Einfache Gruppen.}
\end{center}

Die Sätze, die wir in den vorhergehenden Paragraphen kennen gelernt haben, gestatten ziemlich weitgehende Schlüsse über die Natur der Gruppen. Sie zeigen, dass wenigstens bei den niedrigeren Gradzahlen die metacyklischen (und cyklischen) Gruppen entschieden überwiegen. Um so höheres Interesse beanspruchen
%CONT>>
